\section{Background}
\paragraph{Disk Harmonics.}
The 1D translation bispectrum (which is invariant to group action on the circle $S^1$) is composed of Fourier coefficients, which are constructed using Fourier modes, which are eigenfunctions of the Laplacian on $S^1$. Thus, naturally, if we wish to compose a bispectrum invariant to disk rotations, we should do so with disk harmonic coefficients, which are constructed using disk harmonics, which are the eigenfunctions of the Laplacian on the unit disk $\mathbb{D} = \{ x \in \mathbb{R}^2 : \|x\| < 1 \}$. In polar coordinates \( (r, \theta) \in [0,1) \times [0, 2\pi) \), these disk harmonics are given by
\[
\psi_{nk}(r, \theta) = c_{nk} J_n(\lambda_{nk} r) e^{in\theta}
\]
where \( (n, k) \in \mathbb{Z} \times \mathbb{Z}_{>0} \), \( J_n \) is the $n$-th order Bessel function of the first kind, \( \lambda_{nk} \) is its \( k \)-th root of the $n$-th Bessel function, and \( c_{nk} \) is a normalization constant. Thus, radially, $\psi_{nk}(r, \theta)$ looks like a Bessel function whose value is \( \lambda_{nk} \) on the edge of the disk, and angularly, $\psi_{nk}(r, \theta)$ looks like a complex sinusoid.

\paragraph{Disk Harmonic (DH) Transform.} Let $f : \mathbb{D} \to \mathbb{R}$ be a real-valued square integrable function defined on the unit disk. Disk harmonics form an orthonormal basis for square-integrable functions (like Fourier modes), such that $f$ can be expanded as
$f(r, \theta) = \sum_{n \in \mathbb{Z}} \sum_{k \in \mathbb{Z}_{>0}} a_{nk} \, \psi_{nk}(r, \theta)$,
where $a_{nk} \in \mathbb{C}$ are \textit{disk harmonic (DH) coefficients}. The \textit{disk harmonic (DH) transform} computes these coefficients, expanding $f$ in the disk harmonic basis (akin to the Fourier transform expanding $f$ in the Fourier mode basis). A single coefficient, given by
\[
a_{nk} = \int_0^{2\pi} \int_0^1 f(r, \theta) \, \overline{\psi_{nk}(r, \theta)} \, r \, dr \, d\theta,
\]
encodes the amplitude and phase of $\psi_{nk}$, necessary to reconstruct the signal from the expansion equation above. Importantly, DH coefficients are \textit{equivariant} to rotation of $f$, meaning that a rotation in $f$ will correspond to a rotation in the DH coefficients of $f$.

\begin{proposition}[Equivariance]
More formally, denote $SO(2)$ the group of 2D rotations. Consider $f$ a square-integrable function on the unit disk with DH coefficients $\{a^f_{n, k}\}_{n\in \mathbb{Z}, k\in \mathbb{Z}_{>0}}$. For every $\phi \in SO(2)$, if $f' = f \circ R_\phi$, then for all $n, k \in \mathbb{Z} \times \mathbb{Z}_{>0}$, we have: $a^{f'}_{n, k} = e^{in\phi}\cdot a^{f}_{n, k}.$ Thus, the DH coefficients are equivariant to 2D rotations.
\end{proposition}

In the general case, a function $f$ is \textit{equivariant} to a transformation group $G$ if
\[
f(T_g x) = T_g' f(x), \quad \forall g \in G,
\]
where $T_g$ acts on the input and $T_g'$ is a corresponding action on the output. A function $f$ is \textit{invariant} to $G$ if
\[
f(T_g x) = f(x), \quad \forall g \in G.
\]
\paragraph{Signal Discretization on images.} While the disk harmonics are defined for functions $f : \mathbb{D} \to \mathbb{R}$, practical applications require discretization. We will work with discretized images, i.e., arrays of size $L \times L$ whose pixel values are interpreted as samples of $f$.
Let $p = L^2$ be the total number of pixels such that locations and values of image intensities are denoted as $x_1, \dots, x_p$ and $f_1, \dots, f_p$, respectively. For images of shape \( L \times L \) with spacing \( \Delta x = 2/L \), corresponding to the domain \( [-1,1] \times [-1,1] \), the sampling rate is \( (\Delta x)^{-1} = L/2 \), yielding a Nyquist frequency (i.e., bandlimit) of \( L/4 \). The \textit{Nyquist frequency} is the highest frequency that can be well represented when the signal is sampled at a given rate.

Recall that $\lambda_{nk}$ is the $k$-th root of the $n$-th order Bessel function, and that the disk harmonic $\psi_{nk}$ does not extend radially past $\lambda_{nk}$. Therefore, $\lambda_{nk}$ determines  the number of Bessel function oscillations over the disk and is therefore a quasi-frequency, which is subject to Nyquist criterion. According to the Nyquist criterion, we should only keep the DH coefficients for which $\frac{\lambda_{nk}}{2\pi} \leq \frac{L}{4}$, as higher-frequency components correspond to features beyond the image resolution~\citep{zhao2013fourier}.

\paragraph{Truncation of DH coefficients.} Following closely with \cite{fast_dhcs_2023}, for a given bandlimit $\lambda > 0$, typically $\lambda = \frac{\pi L}{2}$, we keep the DH coefficients for indices $(n,k)$ such that $\lambda_{nk} \leq \lambda$, and order them from lowest frequency to highest $\lambda_0 \leq \cdots \leq \lambda_{m-1}$. The index $j=0, ..., m-1$ represents this order. As such, we now index $\lambda_j$ or equivalently $\lambda_{n_j,k_j}$, where $n_j, k_j$ are the $n,k$ of the frequency at index $j$. Importantly, Bessel function roots do not repeat over order, so we can uniquely identify a disk harmonic via $\lambda_{n_j,k_j}$. As such, we also index $\psi$ by $j$, denoting $\psi_j = \psi_{n_j k_j}$.

Table~\ref{tab:indices} shows the indices $j=0, ..., m-1$ organized in a table labeled by $n, k$ for an image of size $8 \times 8$ and $\lambda = \frac{\pi L}{2}$. Further, for a given bandwidth $\lambda$ that determines a number $m = m_\lambda$, we define the maximum angular frequency (maximum order of the Bessel functions) as \( N_m = \max\{n_j \in \mathbb{Z} : j \in \{0, \ldots, m-1\}\} \) , and the maximum Bessel root as \( K_n = \max\{k \in \mathbb{Z}_{>0} : \lambda_{nk} \leq \lambda\} \) for \( n \in \{-N_m, \ldots, N_m\} \). For example, Table~\ref{tab:indices} has $m=46$, $N_m = 10$, and $K_2 = 4$.
\begin{table*}[t]
\centering
\scriptsize
\resizebox{\textwidth}{!}{
\begin{tabular}{c|ccccccccccccccccccccc}
\hline
$n \rightarrow$ & -10 & -9 & -8 & -7 & -6 & -5 & -4 & -3 & -2 & -1 & 0 & 1 & 2 & 3 & 4 & 5 & 6 & 7 & 8 & 9 & 10 \\
$k \downarrow$ &  &  &  &  &  &  &  &  &  &  &  &  &  &  &  &  &  &  &  &  &  \\
\hline
1 & 44 & 38 & 30 & 25 & 19 & 15 & 10 & 6 & 3 & 1 & 0 & 2 & 4 & 7 & 11 & 16 & 20 & 26 & 31 & 39 & 45 \\
2 & - & - & - & 48 & 40 & 32 & 23 & 17 & 12 & 8 & 5 & 9 & 13 & 18 & 24 & 33 & 41 & 49 & - & - & - \\
3 & - & - & - & - & - & - & 42 & 34 & 27 & 21 & 14 & 22 & 28 & 35 & 43 & - & - & - & - & - & - \\
4 & - & - & - & - & - & - & - & - & 46 & 36 & 29 & 37 & 47 & - & - & - & - & - & - & - & - \\
\hline
\end{tabular}
}
  \caption{Indices $j=0, ..., m-1$, for $m=46$ organized in a table labeled by $n, k$ for an image of size $L\times L = 8 \times 8$ and $\lambda = \frac{\pi L}{2}$. The columns represent the Bessel orders, i.e., the angular frequencies. The rows represent the Bessel root indices, i.e., the radial frequencies.}
  \label{tab:indices}
\end{table*}

\paragraph{Multi-Reference Alignment (MRA)}
In classical (translational) MRA, the goal is to recover a ground truth signal $f$ from multiple noisy and translated observations of the signal $\tilde{f}_i = f + n_i$, for $i=1, \ldots, I$, and where the noise variables $n_i$ are i.i.d. with $n \sim \mathcal{N}(0, \sigma^2)$. Bispectrum MRA does this by mapping all noisy translated observations to their bispectra $b^{\tilde{f}_i}=b_i \in \{b_1, ..., b_I\}$ and computing the average bispectrum $\mathbb{E}[b_i] = b^f + \delta$ , where $b^f$ is the bispectrum of the true signal $f$ and $\delta$ is a bias term introduced by noise propagation. Theoretically, the ground truth signal can be recovered by correcting for the bias term and mapping the corrected bispectrum back to image space.

\cite{bendory2018bispectrum} has derived the bias term for the 1D translation bispectrum so that MRA can be performed on translated 1D signals. In our theory section, we derive a bias term for the 2D disk bispectrum so that MRA can be performed on rotated 2D signals.
